\chapter{Discussion}
\label{chap:discussion}

This chapter discusses the obtained results by summarising the key insights and findings. Furthermore, it provides a critical perspective on the results chapter (\autoref{chap:results}) by examining the results against the study's underlying assumptions and limitations.

\section{Findings and Takeaways}
\label{sec:findings_takeaways}

Main lessons from the staged training procedure, the Phase~2 multi-objective fine-tuning, and the subsequent zero-shot WLTC evaluation are presented while the central message is, that, the proposed approach can learn charge-sustaining, low-\ac{NOx} behaviour in the surrogate environment. This result remains bounded by the selected reward terms and the comparison with the rule-based baseline.

\subsection{Overall Outcome}

% \note{maybe here a: "because of the sheer amount of different experiments, each with a unique purpose in the overall pipeline, a single table that summarises what each experiment was for and what it yielded or sth like dis}

\paragraph{Reinforcement learning as an HEV emission controller.}
The staged curriculum-learning approach can be considered successful within the restrictions of the conducted experiments. The methodology proposed in \autoref{chap:methodology} first isolates the speed-tracking task and then fine-tunes the resulting policies with additional \ac{NOx} and \ac{SOC}-squared penalties. This second stage reduces cycle \ac{NOx} and \ac{SOC} drift by approximately one order of magnitude each while preserving the speed-tracking performance of the two stable agents \ac{PPO} and \ac{SAC} (see \autoref{tab:phase2_headline}).

The main conclusion is therefore not simply that the agents can track a speed trajectory, but that speed tracking, low \ac{NOx}, and \ac{SOC} sustainment can be combined through a scalarised reward and a constrained reward-weight search. By transferring a speed-tracking policy into Phase~2, the multi-objective fine-tuning was able to extent the learned capability without having to relearn basic controllability from scratch.

\paragraph{Charge-sustaining emission reduction.}
The main contribution should be framed as \ac{NOx} reduction under a charge-sustaining constraint. A comparison based on \ac{NOx} alone would be misleading, because Phase~1 already produced an apparently clean seed: SAC seed~8 reached only $2.57$~g cycle \ac{NOx}, but did so by draining the battery. Only one of the 30 Phase~1 seeds was relatively charge-sustaining, but emitted $60.9$~g \ac{NOx}. By contrast, the Phase~2 policies reduce \ac{SOC} drift substantially and improve on this fair Phase~1 comparator by a factor of $4$ to $30$ in cumulative \ac{NOx}.

This distinction is important for interpreting the whole thesis. Low \ac{NOx} is only meaningful if the controller does not obtain it by simply utilising battery energy. The fair result is therefore not "minimum \ac{NOx} at any cost", but low \ac{NOx} while maintaining battery \ac{SOC} within the specified evaluation tolerance.

\subsection{Mechanistic and Algorithmic Insights}

\paragraph{Mechanism of the \ac{NOx} reduction.}
The Phase~2 emission reduction is mechanistically interpretable rather than only a black-box policy effect. It is associated with a shift towards lower \ac{ICE} torque, where the tailpipe \ac{NOx} rate is substantially lower (\autoref{fig:phase2_nox_rate_vs_torque}). Since the Phase~1 policies often overcharged the battery through high-load engine operation, enforcing charge sustainment in Phase~2 removes this wasteful operating mode and reduces \ac{NOx} and fuel consumption simultaneously.

This simultaneous reduction should not be interpreted as a general \ac{NOx}--fuel synergy. It occurs only because the Phase~1 policies were using additional engine work to charge the battery beyond sustainment. Once this behaviour is penalised, both fuel and \ac{NOx} decrease. At the same time, the \ac{BSFC} map alone is not the right lens for judging the \ac{NOx} emission, because the low-\ac{BSFC} region overlaps with high-\ac{NOx} operation. A policy that moves away from the better \ac{BSFC} can therefore still be desirable if the objective is tailpipe \ac{NOx} reduction.

\paragraph{Learning rate as a fine-tuning asset.}
PPO's small learning rate appears to be beneficial during Phase~2 fine-tuning. Under the changed reward function, the small update steps keep PPO close to its Phase~1 speed-tracking policy, which helps avoid catastrophic forgetting and yields the most stable seed population. What looked like a Phase~1 disadvantage, due to slow convergence, therefore becomes useful once the objective changes.

This explanation remains a hypothesis rather than a proven causal mechanism. Learning rate alone cannot explain the full pattern: SAC uses a higher learning rate than TD3 but remains stable, whereas TD3 becomes unstable. The observed stability therefore likely depends on the interaction between learning rate, entropy regularisation, policy stochasticity, and on- versus off-policy updates. A dedicated learning-rate ablation would be required to confirm this interpretation.

\paragraph{TD3 instability.}
\ac{TD3} cannot be recommended for this control problem in its current configuration. If the two divergent runs are excluded, the remaining eight Phase~2 seeds reach a mean \ac{RMSE} of $3.96$~km/h; however, seeds~1 and~3 diverge to \ac{RMSE} values of $6.78$ and $35.57$~km/h. This $20\,\%$ seed-failure rate, combined with the algorithm's overall instability, represents a clear robustness failure, even if some seeds remain competitive with SAC in speed tracking.

A plausible explanation for this behaviour is the effective disabling of delayed actor updates through the Phase~1 hyperparameter selection \texttt{policy\_delay}$=1$. As discussed in \autoref{subsec:theory_td3}, delayed actor updates are one of the defining mechanisms of \ac{TD3}. The Phase~1 speed-tracking objective may have been simple enough to solve without this mechanism, whereas the increased complexity of Phase~2 could have benefited from it. For a controller, the relevant question is whether training reliably produces usable policies; under the present conditions, TD3 does not meet that requirement.

\subsection{Policy Selection and WLTC Generalisation}

\paragraph{Best seed per algorithm.}
Ranking the validation seeds by the composite score on the Phase~2 experiments selects PPO seed~4 ($J = 4.23$), SAC seed~7 ($J = 2.06$), and TD3 seed~9 ($J = 8.67$). Because TD3 is unstable, the practical comparison is between PPO and SAC. SAC seed~7 is the cleanest individual policy across all experiments, whereas PPO provides the most stable seed population. Thus, SAC offers the strongest single candidate for continued development.

\paragraph{Generalisation to WLTC.}
The zero-shot WLTC evaluation gives a more cautious view of the learned policies. PPO and SAC track the WLTC reference more accurately than the rule-based controller under the isoSOC comparison condition (PPO \ac{RMSE} $0.81$~km/h, SAC $1.64$~km/h, rule-based $2.29$~km/h). However, the 10-seed averages emit about $17\,\%$ more \ac{NOx} ($\sim 172$ versus $147$~mg/km) and burn $1.6$ to $2.2\times$ more fuel than the rule-based controller. All seeds exceed the Euro~6 limit of $80$~mg/km except the one near-threshold SAC seed 7.

The WLTC result therefore demonstrate generalisation capabilities but also highlight the persiting importance of a rule-based solution. The learned policies transfer their speed-tracking competence to a cycle that was not used during training, but the population-level emission and fuel performance do not yet surpass the expert-designed rule-based solution.

\paragraph{The fuel and CO\textsubscript{2} blind spot.}
The main reason for the WLTC fuel penalty is that fuel economy was not part of the reward. The Phase~2 fuel reduction on the staircase cycle is only a side effect of charge-sustaining \ac{NOx} optimisation, not evidence of fuel-optimal control. On the WLTC, this limitation becomes visible: PPO keeps the engine on throughout the cycle, and both learned controllers burn more fuel than the rule-based baseline.

The controller should consequently be described as \ac{NOx}- and \ac{SOC}-oriented, not energy-optimal. Adding an explicit fuel or CO\textsubscript{2} term is the most direct next extension of the reward formulation.

\section{Limitations of the Study}
\label{sec:limitations_of_the_study}

\subsection{Scope and Generalisability}
\label{subsec:discussion_scope_generalisability}

\paragraph{Engine-specific models.} First and foremost, the generalisability of this study is limited by the highly customised plant \ac{LSTM}s used to model the EA189 engine. Consequently, if a different power-split \ac{HEV} is to be controlled, the developed models are unlikely to be directly applicable because of differences in internal powertrain functionality. However, although the trained models themselves are not transferable, the underlying methodology remains valid.

\paragraph{Assumption of near-complete plant observability.} Compared with the prior LDCI2027 work, which incorporated only four observations---namely target velocity, actual velocity, \ac{NOx}, and \ac{SOC}---this study uses an extended observation vector with 12 observations. It also includes the SCR1 wall-temperature variable, which is known not to be available during deployment. This creates a simulation-to-reality gap that hinders future deployment. Therefore, each observation should be reassessed to determine whether it is a feasible member of the observation vector.

\paragraph{Single run rule-based baseline.} The comparison with the rule-based solution proposed by Frey et al.\ \cite{frey_data-based_2025} should be interpreted with caution. The learned policies are reported as 10-seed averages to account for stochastic variability during training, whereas the rule-based controller is represented by a single evaluation. This does not invalidate the benchmark, since a rule-based controller does not exhibit seed-dependent training variance under identical conditions, but it makes the variability assessment asymmetric. The comparison is therefore informative as a practical reference point, but it should not be interpreted as a fully statistically balanced superiority test.

\subsection{Experimental Design and Training Procedure}

\paragraph{Limited training trajectory diversity.} The training speed trajectory had several characteristics that could potentially have been modified to accelerate training, improve robustness at speeds above 150~km/h, which were not included in the training range, or produce policies that change speed more reliably and rapidly while remaining efficient. Possible improvements include increasing the maximum speed of the training cycle or shortening the plateaus in the staircase cycle.

\paragraph{Different HPO depths.} Contrary to the initial assumption that only the trials within one cluster \ac{HPO} job would be considered for Phase~1, Optuna considered all entries in \texttt{study\_journal.log}. As a result, the study included 70 PPO trials, 45 TD3 trials, and 44 SAC trials, rather than the intended 20 trials per algorithm. This does not invalidate the resulting models, but it does make the \ac{HPO} depth uneven across algorithms because premature test runs were included in the final optimisation history.

\paragraph{Ordering of analysis steps.} All multi-step experimental pipelines require a decision on how the individual stages are ordered. Instead of the methodology used here, each Phase~1 hyperparameter configuration could have been evaluated across 10 seeds before selecting the best configuration. This would have provided a statistically more robust basis for model selection, but it would also have increased the computational cost substantially. The chosen methodology therefore reflects a compromise between computational feasibility and experimental reliability. Nevertheless, a more exhaustive seed-wise evaluation before checkpoint selection could further improve the robustness of the reported algorithm ranking.

\paragraph{Phase~2 transfer strategy.} In the current experimental design, both the Phase~1 checkpoints and the corresponding hyperparameters are transferred directly to Phase~2. This strategy reduces computational cost and provides stable starting policies that have already learned the basic speed-tracking task. However, it may also constrain the search for control strategies that are better suited to the full multi-objective reward, since the inherited policies were specialised to the speed-tracking-only objective. The same limitation applies to the transferred hyperparameters: an architecture or update schedule that is effective in Phase~1 is not necessarily optimal once \ac{NOx} reduction and \ac{SOC} sustainment are introduced. For example, SAC's three-layer network may have supported its successful adaptation, whereas TD3's Phase~1 setting \texttt{policy\_delay}$=1$ may have been detrimental in Phase~2. Since training from scratch and Phase~2-specific hyperparameter optimisation were not evaluated directly, this transfer trade-off remains hypothetical.

\subsection{Reward and Control Formulation}

\paragraph{Fine-tuning modifications.} Furthermore than the general Phase~2 transfer trategy, two specific adjustments could have further improved the overall quality of the work.

\begin{itemize}
  \item First, the mechanical shaping terms, in particular the mechanical braking penalty defined in \autoref{eq:pmech}, could have been removed after Phase~1 training. These terms were included to guide early learning, because the agents were not able to stabilise speed tracking without a more restrictive reward function, for example by avoiding continuous braking. However, once speed-tracking competence had been established, removing these shaping terms could have given the \ac{DRL} agents more freedom to explore alternative control strategies without this manually imposed restriction.
  \item Second, fine-tuning is often accompanied by a reduction in the learning rate. This was not applied in this study, since all hyperparameters were transferred unchanged from Phase~1 to Phase~2. Although this still produced models that solved the multi-objective control problem successfully, a reduced learning rate may have stabilised TD3 training or enabled further policy refinement.
\end{itemize}

\paragraph{Over-regulation of \ac{SOC} and speed tracking.} The final models are partly a consequence of the acceptance thresholds defined in this study. In particular, the comparison with the rule-based solution highlights the very tight speed-tracking capability of the learned policies. Relaxing the speed-tracking constraint in the composite score could have allowed the optimisation to favour controllers that more closely match the rule-based solution in terms of \ac{NOx} minimisation, at the cost of a small loss in tracking precision, for example on the order of 1--2~km/h. Similarly, relaxing the \ac{SOC} constraint could have enabled greater use of \ac{EM2} and thereby further reduced \ac{NOx}. Overall, this limitation arises from heuristically defined tolerance thresholds for speed error and \ac{SOC} deviation (\autoref{subsec:impl_phase2_pipeline}), which could be aligned more explicitly with the \ac{FVV} project requirements that provide the framework and motivation for this thesis.

\subsection{Scope of the Control Objective}

\paragraph{Component longevity.} The implementation already includes simplified measures to support component durability, such as the engine on/off cooldown and limits on harmful \ac{SOC} states. However, these mechanisms do not constitute a holistic component-longevity model. They capture only selected mechanical-engineering considerations at a simplified level, whereas a more complete durability-aware controller would require closer integration of powertrain-domain expertise.

\paragraph{Single-pollutant objective.} The \emph{most severe limitation} is arguably that, although this study demonstrates successful reductions in \ac{NOx} emissions, the environmental impact of vehicle operation cannot be described by \ac{NOx} alone. Other pollutants and climate-relevant quantities, including the main well-mixed greenhouse gases carbon dioxide (CO\textsubscript{2}), methane (CH\textsubscript{4}), and nitrous oxide (N\textsubscript{2}O), as well as ozone (O\textsubscript{3}), also contribute to the regulatory and environmental assessment of a powertrain \cite{core_writing_team_climate_2023}. Consequently, a controller that optimises only \ac{NOx} may shift the operating strategy in ways that are beneficial for one pollutant but unfavourable for others, especially since N\textsubscript{2}O can be a potent greenhouse-gas by-product of \ac{SCR} aftertreatment. The results should therefore be interpreted as evidence for the feasibility of \ac{NOx}-aware control, not as proof of a globally emissions-optimal strategy. \textbf{A more complete formulation would require a multi-pollutant objective that explicitly accounts for the relevant emission types and their trade-offs.}

\section{Sustainability Analysis}

Following the UPC guidelines \cite{universitat_politecnica_de_catalunya_guia_2023}, this section assesses the sustainability of the thesis along the three classical dimensions --- environmental, economic and social --- and, for each, distinguishes between the \emph{development} of the academic work itself, the \emph{execution} of the project it proposes (here a hypothetical deployment of the learned controller), and the associated \emph{risks and limitations}. Consistent with the proof-of-concept scope defined in \autoref{sec:problem_description}, and because the environmental motivation is the very reason this topic was chosen, the analysis is weighted towards the environmental dimension; the economic and social dimensions and the ethical implications are treated more briefly. In line with the UPC guideline, the emphasis lies on the \emph{ability} to analyse these impacts rather than on the absolute sustainability level of the work.

\subsection{Environmental Dimension}

The environmental footprint of an applied machine-learning thesis is best understood as two opposing contributions: an \emph{intrinsic} cost caused while developing the controller, and an \emph{extrinsic} effect that the controller could deliver if deployed. A meaningful assessment must weigh the former against the latter.

\paragraph{Development of the work - the intrinsic cost.} The dominant environmental cost of this thesis is the electrical energy consumed during training and hyperparameter optimisation. The cumulative computational effort is not negligible. A single seed run of $4{,}000{,}000$ timesteps requires between approximately $8$\,h (PPO) and $15$\,h (SAC) of wall-clock time at $n_{\text{envs}} = 20$ (\autoref{tab:phase1_summary}), and the study comprises tens of such seed runs across both phases, together with a larger number of partially completed hyperparameter-optimisation trials. Several design choices helped to limit this footprint: median pruning terminated unpromising \ac{HPO} trials early instead of running them to completion; the plant model migration to the CTTC-UPC library reduced per-step inference time; environment vectorisation improved hardware utilisation; and the curriculum-inspired learning strategy fine-tuned Phase~2 from the Phase~1 checkpoints rather than training from scratch, which likely reduced the number of timesteps required to reach a proficient policy, as discussed earlier in this chapter. A further mitigating factor is the location of the computation: the training was executed on the Joan Francesc Fernandez (JFF) cluster at CTTC \cite{centre_tecnologic_de_transferencia_de_calor_cttc-upc_heat_2026}, which is located in Spain. At the time of comparison, Spain had a comparatively low electricity carbon intensity (168~gCO$_2$eq/kWh) relative to, for example, Germany (394~gCO$_2$eq/kWh); these snapshot values were taken from \cite{electricity_maps_live_2026}. This illustrates the more general principle that the carbon intensity of \ac{AI} development can be lowered simply by scheduling compute in regions with a high share of renewable energy. A precise $\mathrm{CO}_2$e estimate would require the cluster's energy accounting and is therefore not quantified here. Importantly, however, this footprint represents a one-time offline cost.

\paragraph{Hypothetical project execution - the extrinsic potential.} The environmental relevance of the proposed controller is that the resulting artefact is \emph{software}: a trained policy that executes in under a millisecond per step on existing electronic control hardware and requires no additional material or manufacturing. If such a controller delivered a robust per-vehicle reduction in tailpipe emissions, the one-off offline training cost could be amortised over the operating life of every vehicle using it, and even moderate improvements could become relevant at higher scale. The extrinsic sustainability contribution of this thesis should therefore be understood primarily as deployment potential and research insight rather than as a directly quantified emission reduction. The study shows how a scalable software controller can learn emission-aware behaviour under selected conditions, and it documents which design choices, constraints and failure modes shape that behaviour. These findings cannot yet be converted into a fleet-level environmental benefit, because their future use within the \ac{LDCI2027} group and beyond is uncertain. Nevertheless, they remain relevant because technological progress often depends on such intermediate evidence, negative results and design guidance.

\paragraph{Environmental risks and limitations.} Three caveats accompany this potential: First, the objective optimises \ac{NOx} alone. As discussed earlier in this chapter, a controller that minimises \ac{NOx} may shift other quantities (carbon monoxide, hydrocarbons, particulates, $\mathrm{CO}_2$) unfavourably, so the demonstrated reduction is not proof of a globally emissions-optimal strategy. Second, the comparison with the rule-based reference on the realistic \ac{WLTC} limits the strength of the sustainability claim (\autoref{sec:wltc-comparison}). The learned policies track the target more tightly, but their 10-seed averages consume $1.6$--$2.2\times$ the fuel and emit about $17\%$ more \ac{NOx}; only the selected SAC seed~7 reaches a clearly favourable \ac{NOx} result, and it still does so with higher fuel consumption. A net environmental benefit over the established rule-based solution is therefore not yet demonstrated. Third, the entire evaluation rests on \ac{LSTM} plant surrogates rather than a physical vehicle, so simulation-to-reality effects remain unresolved. The honest reading is therefore that the work demonstrates the \emph{feasibility} of emission-aware learned control, not a delivered environmental gain.

\subsection{Economic Dimension}

\paragraph{Development of the work - the development cost.} Expressed as the invoice a client would receive, the cost of the work is dominated by the author's time, including literature review, study and system design, implementation and debugging, and academic writing. The project started in January 2026 and extended over approximately 21 weeks until June 2026. Assuming an average workload of 30 hours per week yields 630 hours of work; at a hypothetical rate of EUR~20 per hour, this corresponds to EUR~12{,}600. These values are heuristic, so it should be noted that the estimated cost would increase, for example, if a German salary level were assumed, which is typically higher \cite{statistisches_bundesamt_earnings_nodate}. Conversely, EUR~20 per hour is already a comparatively generous assumption for student work in Catalonia.

All software (Stable-Baselines3 \cite{raffin_stable-baselines3_2021}, Optuna \cite{akiba_optuna_2019}, the underlying runtime of the CTTC-UPC models, and the Python scientific stack) is open-source, so the licensing cost is zero, and the marginal monetary cost of computation on an existing academic cluster is small and, crucially, not borne by the author. However, a two-month Claude Pro \cite{anthropic_claude_nodate} subscription was purchased for support during development of the thesis at EUR~21.42 per month. The hardware used for the work, namely the author's personal laptop, had already been purchased two years earlier and is therefore not included in this calculation. Overall, this leads to a total estimated cost of \emph{EUR~12{,}642.84}.

\paragraph{Hypothetical project execution - the project viability.} A learned \ac{EMS} is economically attractive because, once trained, it is software with a near-zero marginal cost per vehicle. It could also replace part of the meticulous manual calibration required for rule-based controllers (\autoref{sec:problem_description}) with an automated, data-driven procedure, potentially reducing the engineering effort needed to develop and recalibrate an \ac{EMS} across platforms and operating conditions. Avoiding regulatory non-compliance penalties represents a further indirect economic benefit.

\paragraph{Economic risks and limitations.} The primary recurring cost arises from the engine-specific nature of the plant surrogates (\autoref{subsec:discussion_scope_generalisability}): a different powertrain would require new data collection, model training and controller re-optimisation, which eliminates the scalability advantage. In addition, the simulation-to-reality gap implies substantial validation costs before deployment. The approval of a learned, less interpretable controller for a safety-relevant function may also be expensive or uncertain. In such a scenario, a simpler rule-based controller might remain the more economical choice.

\subsection{Social Dimension}

\paragraph{Development of the work - the impact on the developers of the work.} On a personal level, the thesis required the author to integrate several previously separate areas: reinforcement learning, automotive powertrain engineering, and sustainability assessment. The work was motivated by the aspiration to direct technical skills towards an environmental problem, which made it personally fulfilling. It was carried out within the international \ac{LDCI2027} collaboration in accordance with good professional and academic practice, and uses inclusive, non-discriminatory language throughout the report. The work complies with the UPC Code of Ethics \cite{universitat_politecnica_de_catalunya_code_2022}.

\paragraph{Hypothetical project execution - the social impact of the project.} The intended positive effects are diffuse but relevant. Lower tailpipe \ac{NOx} can benefit public health and urban air quality, since \ac{NOx} emissions are associated with respiratory disease \cite{european_union_regulation_2024, european_union_regulation_2007, core_writing_team_climate_2023}. Vehicle manufacturers and regulators may also benefit from methods that support emission compliance. More broadly, the work contributes to the societal effort to reduce the environmental impact of mobility, which is connected to the wider challenge of mitigating global warming and its associated social and health consequences \cite{core_writing_team_climate_2023}.

\paragraph{Social risks and limitations.} A learned controller is less transparent than an explicit rule set, which makes it harder for engineers, regulators and users to understand why a specific action was taken. Beyond this interpretability concern, the main social risks are:

\begin{itemize}
  \item \textbf{Safety and robustness.} If the policy behaves unexpectedly outside its validated operating domain, the consequences could affect drivers, passengers and other road users, even though the controller is not intended as an autonomous-driving system (\autoref{sec:objectives}).
  \item \textbf{Human oversight and deskilling.} Excessive reliance on learned controllers could weaken manual-calibration expertise and reduce the ability of engineers to intervene when model behaviour malfunctions. This represents a standard risk of data-driven methods and is not specific to this study only.
\end{itemize}

Regarding data-protection concerns, the work only uses simulated engine data and no personal data, so the GDPR is not engaged beyond the standard institutional commitments.

\subsection{Ethical Implications}
\label{subsec:ethics}

Reflecting on the consequences of this study raises several ethical considerations that are specific to a learned, sustainability-motivated controller.

\paragraph{Conflict of interest and research independence.} This thesis is embedded in the \ac{LDCI2027} project under the \ac{FVV} \cite{fvv_e_v_-_forschungsvereinigung_verbrennungskraftmaschinen_fvv_nodate}, an industry research association for \emph{combustion engines}. Work funded to improve combustion-based powertrains inherits a structural incentive to present that pathway favourably and, indirectly, to prolong the relevance of the \ac{ICE}. This potential conflict is disclosed here in the interest of transparency. The analysis throughout the thesis, including the unfavourable \ac{WLTC} comparison in \autoref{sec:wltc-comparison}, is reported independently of that incentive, and no result has been framed to support the sponsor's technology.

\paragraph{Reward design as a choice.} A reinforcement-learning controller optimises exactly what its reward function encodes. The relative weighting of emissions, speed tracking and charge sustainment ($W_{\text{emi}}$, $W_{\text{soc}^2}$ and the shaping terms) is therefore not a neutral engineering parameter but an implicit ethical decision about whose interests are prioritised: public health (minimal \ac{NOx}), driver convenience (tight tracking), or manufacturer cost (fuel and component longevity). Making these weights and their selection procedure explicit (\autoref{subsec:method_phase2}) is an ethical safeguard, because it allows the interests encoded in the objective to be inspected rather than hidden inside a tuning constant.

\paragraph{Accountability for an opaque controller.} Because the policy is a black box acting on a safety-relevant powertrain function, a malfunction creates a responsibility gap: it must not be possible to displace blame onto "the algorithm". Responsibility must remain with identifiable human and institutional actors, including the engineers who design the reward and validation procedure and the organisation that eventually deploys the controller. This is the ethical motivation for retaining human oversight and a verifiable validation process before any real-world use.

\paragraph{Precautionary deployment.} The controller is validated only against \ac{LSTM} surrogates and a synthetic cycle, and the objectives of this thesis explicitly exclude autonomous driving or production use (\autoref{sec:objectives}). Treating the result as a proof of concept rather than a deployable artefact is therefore the ethically required precautionary position. An unvalidated learned controller should not influence a physical vehicle on public roads until its real-world behaviour and failure modes are demonstrably safe.

\paragraph{Greenwashing.} Given the framing of this work, a positive proof-of-concept result must not be overstated. Presenting selected, simulation-only \ac{NOx}-aware behaviour as evidence of a "clean" or "sustainable" powertrain would constitute greenwashing. Similarly, overstating improvements in a combustion-based \ac{HEV} could lend overstated legitimacy to continued combustion-engine development and thereby slow the broader transition to fully electric mobility or other low-carbon alternatives. Any such claim must therefore be accompanied by the limited scope of the work: a single-pollutant, in-simulation feasibility study. The reporting of unfavourable outcomes --- including the \ac{WLTC} fuel penalty, the single-pollutant scope, and the \ac{TD3} instability --- follows directly from this commitment.

\subsection{Relationship with the Sustainable Development Goals}

The thesis relates most directly to the following United Nations Sustainable Development Goals (SDGs) \cite{united_nations_17_nodate}:

\begin{itemize}
  \item \textbf{SDG~3 (Good Health and Well-being)}, target~3.9: reducing tailpipe \ac{NOx} directly addresses illness and deaths from ambient air pollution.
  \item \textbf{SDG~9 (Industry, Innovation and Infrastructure)}, target~9.4: the work exemplifies upgrading existing industry for greater resource-use efficiency and cleaner technology through innovation.
  \item \textbf{SDG~11 (Sustainable Cities and Communities)}, target~11.6: lower vehicle emissions reduce the adverse per-capita environmental impact of cities, chiefly through improved air quality.
  \item \textbf{SDG~12 (Responsible Consumption and Production)}: more efficient powertrain operation makes more economical use of fuel per kilometre travelled.
  \item \textbf{SDG~13 (Climate Action)}: approaching fuel and emission reduction supports the broader decarbonisation effort.
\end{itemize}

The contribution to each goal is indirect and conditional on eventual deployment, but the intent and direction of the work are clearly aligned with their spirit.
